En este apartado se analiza la factibilidad legal del proyecto, donde se verifica el cumplimiento de los requisitos jurídicos y si esta dentro de las leyes, normas y reglamentos vigentes del marco jurídico sin llegar a infringir leyes, derechos de autor o licencias de terceros.  

La revisión de factibilidad legal del proyecto se estructura en los siguientes campos de análisis: 

\begin{itemize}
    \item Protección de datos personales y privacidad.
    \item Licenciamiento de las herramientas tecnológicas para el proyecto.
    \item Propiedad intelectual y derechos de autor
    \item Cumplimiento de Términos de Servicio (ToS) y uso de datos externos.
    \item Términos y condiciones de uso. 
\end{itemize}

\subsubsection{Protección de datos personales y privacidad}

El sistema gestiona la información de los perfiles de usuario que forman parte del ecosistema del prototipo de aplicación móvil por lo que dada la naturaleza de este hecho es necesario procesas datos personales de los usuarios al igual que datos relacionados con la geolocalización de toda la área de acción que comprende el sistema. Se trabaja con información sensible incluyendo direcciones de correo electrónico, credenciales de acceso y coordenadas de geolocalización en tiempo real, por consiguiente, la operación del sistema se rije bajo la \textbf{Ley de Federal de Protección de Datos Personales en Posesión de los Particulares} o LFPDPPP por sus siglas en español, lo que garantiza el tratamiento legítimo, controlado, seguro e informado de la información de los usuarios. 

Todos los artículos y preceptos jurídicos analizados en este apartado tienen sus fundamento estricto en la Ley Federal de Protección de Datos Personales en Posesión de los Particulares (LFPDPPP) vigente \cite{LFPDPPP_2025}. 

Para asegurar el cumplimiento de los principios de la ley LFPDPPP y las reglas de negocio del sistema, se implementan la siguientes normativas dentro de la tabla \ref{tab:LFPDPPP}:

\textbf{Nota}: Es importante mencionar que la tabla \ref{tab:LFPDPPP} hace referencia a reglas y requerimientos previamente establecidos en el capitulo \ref{cap:analisis_sistema}, específicamente en las secciones \ref{sec:requerimientos} y \ref{sec:reglas_negocio}.

\begin{table}[H]
\centering
\renewcommand{\arraystretch}{1.3} 
\begin{tabular}{|>{\raggedright\arraybackslash}p{0.25\textwidth}|>{\raggedright\arraybackslash}p{0.68\textwidth}|}
\hline
\textbf{\centering Normativa / Artículos LFPDPPP} & \textbf{\centering Descripción y Cumplimiento} \\ \hline

\textbf{Consentimiento expreso e Información}\newline (Artículos 7, 14 y 16) & Conforme al requerimiento RF-46 y la regla RN-16, el sistema hace uso de un Aviso de Privacidad dentro de la aplicación en el acceso inicial de los usuarios, registrando el consentimiento explícito del usuario antes de crear la cuenta. \\ \hline

\textbf{Proporcionalidad en Geolocalización}\newline (Artículos 11 y 12) & En la regla RN-16 se establece el uso exclusivo de la ubicación GPS para funciones que la requieran, como el cálculo de rutas o navegación en vivo. La recolección de datos GPS se detiene una vez terminadas las tareas dentro de la aplicación, sin ceder dicha información a terceros o ejecutar rastreo en segundo plano sin autorización. \\ \hline
\textbf{Medidas de Seguridad de la Información}\newline (Artículo 18) & El requerimiento RNF-14 establece que el almacenamiento de contraseñas se lleva a cabo mediante funciones hash con sal (bcrypt o Argon2). Además, la transmisión de datos opera obligatoriamente bajo el protocolo HTTPS con cifrado TLS (RNF-15), garantizando la protección técnica contra el acceso no autorizado. \\ \hline

\textbf{Derechos ARCO y Disociación de Datos}\newline (Art. 2 fracc. IX, Art. 21 y 24) & El requerimiento RF-11 permite al usuario solicitar la cancelación definitiva de su cuenta. La regla RN-15 establece que los datos personales asociados serán destruidos en un plazo máximo de 7 días. Los reportes ciudadanos son sometidos a un proceso de disociación (anonimización), conservándolos únicamente con fines estadísticos. \\ \hline

\textbf{Transferencia Internacional de Datos}\newline (Art. 2 fracc. XX y Art. 35) & Al utilizar Supabase como infraestructura Backend, la información puede almacenarse en servidores fuera del territorio nacional. Esta condición se estipula explícitamente en el Aviso de Privacidad, respaldando la transferencia en certificaciones internacionales de seguridad (SOC 2) del proveedor. \\ \hline
\end{tabular}
\caption{Cumplimiento de normativas de la LFDPPP y su relación con requerimientos del sistema}
\label{tab:LFPDPPP}
\end{table}

\subsubsection{Licenciamiento de las herramientas tecnológicas para el proyecto.}

El desarrollo del prototipo integra tecnologías, lenguajes y plataformas de base de datos que operan bajo licencias de código abierto, lo que hace al proyecto completamente factible y libre de responsabilidades por infracción de patentes de software comercial.

Todas las herramientas tecnológicas contempladas para el desarrollo del proyecto se establecieron con anterioridad en la tabla \ref{tab:herramientas_proyecto}. La factibilidad legal para el desarrollo del software recae en la utilización de herramientas de código abierto con licencia gratuita, comprobado en la tabla de costos \ref{tab:recursos_tecnologicos}, donde cada una de las herramientas contempla un costo de 0 pesos (considerando el uso de licencia gratuita), lo que mitiga los riesgos por demandas de propiedad intelectual o uso indebido de software privado.    

El uso de software de pago y la compra de licencias se contempla dentro de la sección \ref{sec:factibilidad_economica} cumpliendo con los marcos legales sin incurrir en practicas ilícitas como el uso de licencias ilegales. La unica excepción recae en el uso de APIs de operación de la aplicación (Google Street View Static API y Qwen3.7-Flash) ya que el consumo  sale de los limites de las cuotas gratuitas establecidas, por lo que se contempla el pago por uso una vez superada la tarifa estándar.  

El costo de uso de servicios tecnológicos adicionales como Overleaf y Udemy se contempla en la tabla \ref{tab:servicios_adicionales}, donde ya se consideran los costos necesarios para evitar el uso de software de forma ilegal. 

\begin{table}[H]
\centering
\small 
\renewcommand{\arraystretch}{1.3} 
\begin{tabular}{|>{\raggedright\arraybackslash}p{0.35\textwidth}|>{\raggedright\arraybackslash}p{0.45\textwidth}|}
\hline
\textbf{\centering Herramienta} & \textbf{\centering Licencia / Pago / Suscripción} \\ \hline
Visual Studio Code & Gratuito  \\ \hline
Flutter / Dart & Gratuito  \\ \hline
Supabase & Gratuito  \\ \hline
Git / Github & Gratuito  \\ \hline
Google Meet & Gratuito  \\ \hline
Notion & Gratuito  \\ \hline
OpenStreetMap (OSM) & Gratuito  \\ \hline
Google Street View Static API & Pago (al exceder la tarifa estándar) \\ \hline
Qwen3.7-Flash & Pago (al exceder la tarifa estándar) \\ \hline
Maps SDK (para móvil) & Gratuito \\ \hline
Places API (new)  & Gratuito \\ \hline
Registro de desarrollador Google Play & Pago \\ \hline
Overleaf & Suscripción \\ \hline
Udemy & Suscripción \\ \hline
\end{tabular}
\caption{Resumen de licencias, pagos y suscripciones del stack tecnológico.}
\label{tab:licencias_fact}
\end{table}

\subsubsection{Propiedad intelectual y derechos de autor}

Al ser un proyecto desarrollado en el marco académico de la Escuela Superior de Cómputo del Instituto Politécnico Nacional (ESCOM-IPN), la titularidad de los derechos patrimoniales y morales está sujeta al Reglamento Orgánico del IPN.

Acorde con la Ley Federal del Derecho de Autor (LFDA) vigente \cite{LFDA_ley}, y de conformidad con lo establecido en su Artículo 18 y Artículo 19, los derechos morales (perpetuos, inalienables, irrenunciables, imprescriptibles e inembargables) sobre el código fuente, la lógica algorítmica y el diseño de experiencia de usuario correspondena los alumnos desarrolladores: Sánchez Gómez Alan Ivan, Hernández Simón Rebeca y Mendoza Segura Fernando. Se reconoce legal y académicamente la dirección y asesoría técnica de la Dra. Claudia Jisela Dorantes Villa y el Mtro. José Antonio Ortiz Ramírez. Adicionalmente, conforme a los supuestos de protección contemplados en el Artículo 13 de la misma Ley, el código fuente derivado del Trabajo Terminal 2027-A215 es susceptible de ser registrado ante el Instituto Nacional del Derecho de Autor (INDAUTOR) bajo la rama de Programas de Cómputo.

Por otra parte, al ser un Trabajo Terminal desarrollado como requisito académico de titulación dentro de las instalaciones y programas de la ESCOM, el ejercicio de los derechos de explotación patrimonial regulados por el Artículo 27 de la LFDA (los cuales facultan la reproducción, publicación o edición material y electrónica), así como el licenciamiento comercial final del software, quedarán sujetos a los convenios de cesión (en concordancia con el Artículo 30 de la Ley) o titularidad estipulados en el Reglamento Orgánico del IPN.

\subsubsection{Cumplimiento de Términos de Servicio (ToS) y uso de datos externos.}
Este apartado hace referencia al cumplimento de reglas establecidas al consumir APIs y servicios de terceros para no infringir la obligación legal hacia los proveedores de las tecnologías.  

El funcionamiento del prototipo de aplicación móvil depende del consumo de interfaces de programación de aplicaciones (APIs) y bases de datos espaciales provistas por terceros. Para garantizar la factibilidad legal y operativa del sistema, el diseño de la arquitectura tecnológica se desarrolla conforme el estricto cumplimiento de Términos de Servicio (ToS) y las licencias de uso de datos externos de cada proveedor, detallados a continuación: 

\textbf{Ecosistema de Goolgle Maps Platform}: El sistema integra 3 servicios de Google fundamentales para el funcionamiento: Street View Static API, Maps SDK para móvil y Places API. Conforme a los Términos de Servicio de Google Maps Platform, el proyecto cumple con los siguientes lineamientos: 

\begin{itemize}
    \item\textbf{Prohibición de extracción masiva (Scrapping)}: Las 35,100 imágenes extraídas mediante Street View Static API son utilizadas temporalmente para calcular las condiciones de riesgo. El sistema no almacena las imágenes obtenidas del API permanentemente ni busca crear una base de datos que compita con Google, unicamente conserva información útil para el procesamiento del análisis de riesgo. 
    \item\textbf{Uso comercial y cuotas}: El consumo de APIs de operación (Maps SDK y Places API) se mantiene dentro de los lineamientos de la cuota gratuita mensual. Para la parte del preprocesamiento masivo de imágenes, se ha contemplado con anterioridad el pago correspondiente por exceder la tarifa estándar, lo que garantiza que el uso del servicio sea legitimo sin incurrir en practicas ilegales como la evasión de facturación. 
\end{itemize}

\textbf{Modelo de Lenguaje Visual Qwen3.7-Flash}: El análisis de las características del entorno físico se realiza mediante el modelo Qwen3.7-Flash. El cumplimiento de sus Términos de Servicio se aborda mediante el pago por consumo de tokens (estimado en 70.2 millones de tokens de entrada y 17.55 millones de salida). Cabe mencionar que el sistema no infringe las normativas de la API sobre el procesamiento de contenido privado o material protegido por derechos de autor ya que solo procesa imágenes de fachadas y calles públicas. 

\textbf{Cartografía base y OpenStreetMap (OSM)}: Debido a que el sistema implementa el uso de PostGIS y pgRouting para calcular las rutas y guardar los mapas, la red vial es capaz de ser alimentada de datos gratuitos de OpenStreetMap. Para cumplir con la Licencia Abierta de Bases de Datos (ODbL, por sus siglas en inglés) que rige a OSM, la aplicación móvil debe incluir en su interfaz gráfica el reconocimiento y atribución explícita indicando "© Colaboradores de OpenStreetMap", garantizando así el respeto a la propiedad intelectual comunitaria.

\textbf{Infraestructura Backend-BD (Supabase}: La gestión de la base de datos y el almacenamiento en la nube se realiza a través de Supabase bajo su plan gratuito. El proyecto se adhiere a los limites establecidos dentro del plan gratuito asegurando que el volumen de almacenamiento y las transferencias mensuales no excedas los límites previstos. Además, se garantiza que la infraestructura no será utilizada para alojar contenido ilícito o realizar ataques de denegación de servicio.   


\subsubsection{Términos y condiciones de uso. }
Los Términos y condiciones de uso constituyen el contrato que vincula a los desarrolladores del sistema con el usuario final de la aplicación. Su aceptación explicita es un requisito indispensable para el uso de la aplicación ya que es dentro de este apartado donde se estipulan las reglas y condiciones de utilizar el sistema, al igual que la consecuencias de aceptar dichas condiciones, por lo que los usuarios deben tomar expresa responsabilidad en el uso del software.  

El \textbf{Prototipo de aplicación móvil para la generación de rutas peatonales seguras dentro de la alcaldía Gustavo A. Madero} es una herramienta puramente informativa y probabilística, basada en datos estadísticos y ambientales. No garantiza la seguridad absoluta del usuario, ya que el riesgo de un incidente siempre está presente en la vía pública. Por ello, el equipo de desarrollo no se hace responsable por cualquier daño a la integridad física que el usuario pueda sufrir al utilizar el sistema. 

Para garantizar el funcionamiento adecuado de la aplicación y delimitar el alcance de la responsabilidad de los desarrolladores y del Instituto Politécnico Nacional, es indispensable establecer un marco normativo interno que regule la interacción entre el usuario final y la plataforma. este conjunto de normativas, es formalizado a través de los términos y condiciones de uso. Asimismo, este apartado establece los lineamientos estrictos para el uso ético de las funciones críticas del sistema, tales como la emisión de reportes mediante crowdmapping, la compartición de ubicación en tiempo real y la activación del botón de emergencia.

A continuación, en la Tabla \ref{tab:terminos_condiciones} se desglosan y detallan a fondo las cláusulas normativas principales que conforman este acuerdo, especificando la descripción de cada condición jurídica y la responsabilidad operativa directa que asume el usuario al interactuar con el prototipo.

\begin{table}[H]
\centering
\scriptsize
\renewcommand{\arraystretch}{1.4} 
\begin{tabular}{|>{\raggedright\arraybackslash}p{0.22\textwidth}|>{\raggedright\arraybackslash}p{0.38\textwidth}|>{\raggedright\arraybackslash}p{0.32\textwidth}|}
\hline
\rowcolor{gray!15}
\textbf{\centering Tema / Cláusula} & \textbf{\centering Descripción de la Condición} & \textbf{\centering Responsabilidad y Aceptación del Usuario} \\ \hline

\textbf{1. Naturaleza del Sistema y Descargo de Responsabilidad} & 
El sistema provee rutas peatonales sugeridas basándose en un modelo algorítmico de riesgo socio-espacial y reportes ciudadanos. No garantiza la seguridad absoluta ni la evasión total de incidentes delictivos. & 
El usuario comprende que transita por la vía pública bajo su propio riesgo. Libera a los desarrolladores, asesores y al Instituto Politécnico Nacional (IPN) de cualquier responsabilidad civil, penal o administrativa ante cualquier siniestro o perjuicio. \\ \hline

\textbf{2. Autenticación y Credenciales de Acceso} & 
El acceso a las funciones de ruteo y reporte requiere de una cuenta verificada en el sistema (salvo autenticación por terceros como Google). & 
El usuario es el único responsable de salvaguardar la confidencialidad de su contraseña y evitar el uso compartido o no autorizado de su cuenta. \\ \hline

\textbf{3. Privacidad y Datos de Geolocalización} & 
El sistema recopila datos de ubicación (GPS) en tiempo real estrictamente para el cálculo de rutas, navegación, alertas de proximidad y enlaces de seguimiento, conforme al Aviso de Privacidad y la LFPDPPP. & 
El usuario otorga su consentimiento explícito para la recolección transitoria de su ubicación y acepta que el correcto funcionamiento de la aplicación depende de este permiso. \\ \hline

\textbf{4. Sistema de Reportes (Crowdmapping)} & 
La plataforma permite el registro de reportes sobre el entorno físico. La validación de los mismos opera mediante un sistema de reputación de usuarios (escala de 0 a 5 estrellas). & 
El usuario se compromete a emitir información verídica. Acepta que la emisión de reportes falsos o malintencionados derivará en la penalización de su reputación y posible suspensión de la cuenta. \\ \hline

\textbf{5. Uso de Funciones de Emergencia} & 
La aplicación proporciona un botón de pánico que enlaza de manera directa mediante llamada telefónica a los servicios de emergencia oficiales (911). & 
El usuario se compromete a utilizar esta función de manera exclusiva ante emergencias reales, asumiendo las consecuencias legales establecidas por la ley ante falsas alarmas. \\ \hline

\textbf{6. Enlaces de Seguimiento Compartido} & 
El sistema permite la generación de URLs temporales para compartir el trayecto en vivo con terceros. El enlace se invalida automáticamente al finalizar o cancelar la ruta. & 
El usuario asume la total responsabilidad sobre las personas o plataformas con las que decide compartir su enlace de seguimiento en vivo. \\ \hline

\textbf{7. Disponibilidad del Servicio de Ruteo} & 
La generación de rutas seguras depende de la disponibilidad de APIs de terceros (Google, Supabase) y de la conectividad a internet del dispositivo móvil. & 
El usuario acepta que el servicio puede presentar interrupciones, demoras o bloqueos por fallas de conectividad externas al control de los desarrolladores del sistema. \\ \hline

\textbf{8. Modificación y Cancelación de Cuenta} & 
El sistema retendrá los datos del usuario mientras la cuenta esté activa. Ante la eliminación de la misma, los datos personales se destruirán en un plazo máximo de 7 días. & 
El usuario tiene el derecho de cancelar su cuenta en cualquier momento, aceptando que sus reportes previos se conservarán de forma anonimizada para fines estadísticos. \\ \hline

\end{tabular}
\caption{Detalle de los términos y condiciones de uso para el usuario final del sistema.}
\label{tab:terminos_condiciones}
\end{table}

